\section{SQL}
SQL steht für \emph{Structured Query Language} und ist die Standardsprache, um relationale Datenbanken zu definieren, zu füllen, abzufragen und zu verwalten. SQL ist deklarativ: Man beschreibt, \emph{was} man haben möchte, und nicht, \emph{wie} die Datenbank das Ergebnis beschaffen soll. Die Sprache wird in vier Teilsprachen unterteilt; die Transaktionssteuerung (TCL) wird oft als fünfte Teilsprache ergänzt.

\begin{stdtable}{l X X}{Teilsprache & Aufgabe & Wichtige Befehle}
  \icode{DDL} -- Data Definition Language & Schema der Datenbank definieren und verändern & \icode{CREATE}, \icode{ALTER}, \icode{DROP}, \icode{TRUNCATE} \\
  \icode{DML} -- Data Manipulation Language & Daten in bestehenden Tabellen einfügen, ändern und löschen & \icode{INSERT}, \icode{UPDATE}, \icode{DELETE}, \icode{REPLACE} \\
  \icode{DQL} -- Data Query Language & Daten aus einer oder mehreren Tabellen lesen & \icode{SELECT} \\
  \icode{DCL} -- Data Control Language & Zugriff und Rechte für Nutzer steuern & \icode{GRANT}, \icode{REVOKE} \\
  \icode{TCL} -- Transaction Control Language & Transaktionen steuern & \icode{COMMIT}, \icode{ROLLBACK} \\
\end{stdtable}

\subsection{DDL -- Data Definition Language}
Die DDL umfasst alle Befehle, die die \emph{Struktur} einer Datenbank definieren und verändern. Es geht nicht um die Daten selbst, sondern um das Schema. Fast jedes Datenbankobjekt kann mit \icode{CREATE} erzeugt, mit \icode{ALTER} verändert, mit \icode{RENAME}/\icode{CHANGE} umbenannt und mit \icode{DROP} gelöscht werden.

\subsubsection{Datenbank und Tabelle anlegen}
Mit \icode{IF NOT EXISTS} wird eine Exception abgefangen, falls die Datenbank bereits existiert.

\begin{syntaxbox}[SQL]
CREATE DATABASE IF NOT EXISTS kurs_db;

CREATE TABLE Personal (
    PersonalNummer  INT UNSIGNED,             -- kein NOT NULL noetig (ist PK)
    Vorname         VARCHAR(60)  NOT NULL,
    Nachname        VARCHAR(60)  NOT NULL,
    GeburtsDatum    DATE         NOT NULL,
    AbteilungsID    SMALLINT UNSIGNED NOT NULL,
    PRIMARY KEY (PersonalNummer),
    FOREIGN KEY (AbteilungsID) REFERENCES Abteilung(AbteilungsID)
);
\end{syntaxbox}

\subsubsection{Datentypen}
\begin{stdtable}{l X}{Art & Typen}
  Text & \icode{CHAR(n)} feste Länge \(\cdot\) \icode{VARCHAR(n)} variable Länge, maximal $n$ Zeichen \(\cdot\) \icode{TEXT}/\icode{BLOB} sehr große Inhalte \\
  Ganze Zahlen & \icode{INT}/\icode{INTEGER}, \icode{SMALLINT}, \icode{TINYINT} \\
  Gleitkommazahlen & \icode{FLOAT}, \icode{DOUBLE} -- z.\,B. \icode{FLOAT(8,2)} mit 2 Nachkommastellen \\
  Logisch & \icode{BOOL} (intern ein \icode{TINYINT}) \\
  Zeiten & \icode{DATE} \(\cdot\) \icode{TIME} \(\cdot\) \icode{DATETIME} \(\cdot\) \icode{TIMESTAMP} (setzt bei \icode{INSERT}/\icode{UPDATE} automatisch Datum und Uhrzeit) \(\cdot\) \icode{YEAR} \\
\end{stdtable}

\subsubsection{Spaltenattribute und Constraints}
\begin{stdtable}{l X}{Schlüsselwort & Bedeutung}
  \icode{PRIMARY KEY} & Primärschlüssel: eindeutig, nie \icode{NULL}, identifiziert jeden Datensatz. Kann auch aus mehreren Spalten bestehen. \\
  \icode{FOREIGN KEY} & Fremdschlüssel: verweist auf den Primärschlüssel einer anderen Tabelle. \\
  \icode{NOT NULL} & Die Spalte darf keinen \icode{NULL}-Wert enthalten. \\
  \icode{UNIQUE} & Alle Werte der Spalte müssen eindeutig sein. \\
  \icode{DEFAULT} & Standardwert, falls beim Einfügen kein Wert angegeben wird. \\
  \icode{AUTO\_INCREMENT} & Nummeriert Datensätze automatisch aufsteigend ab 1; nur für ganzzahlige, eindeutige Felder (\icode{PRIMARY KEY} oder \icode{UNIQUE}). \\
  \icode{UNSIGNED} & Nur positive Zahlen erlaubt (kein Vorzeichen). \\
\end{stdtable}

\begin{warnbox}[Stolperfalle]
In SQL kann man eine Tabelle auch ohne \icode{PRIMARY KEY} erstellen -- das sollte man aber nicht tun und den Schlüssel bereits bei der Modellierung festlegen.
\end{warnbox}

\subsubsection{Struktur ändern mit ALTER TABLE}
\icode{ALTER TABLE} verändert die Struktur einer bestehenden Tabelle, ohne die Daten zu löschen.

\begin{syntaxbox}[SQL]
ALTER TABLE Personal ADD Gender CHAR AFTER GeburtsDatum;   -- Spalte hinzufuegen
ALTER TABLE Personal MODIFY GeburtsDatum DATETIME;          -- Datentyp aendern
ALTER TABLE Personal CHANGE GeburtsDatum Geburtsdatum DATE; -- Spalte umbenennen
ALTER TABLE Personal DROP Gender;                           -- Spalte loeschen
ALTER TABLE Personal RENAME TO Mitarbeiter;                 -- Tabelle umbenennen

-- Foreign Key nachtraeglich setzen
ALTER TABLE Personal
    ADD FOREIGN KEY (AbteilungsID) REFERENCES Abteilung(AbteilungsID);
\end{syntaxbox}

\subsubsection{Referenzielle Integrität}
Die Tabelle mit dem Primärschlüssel ist der \emph{Parent} (Master), die Tabelle mit dem Fremdschlüssel das \emph{Child}. Die referenzielle Integrität verhindert verwaiste Daten -- etwa eine Bestellung ohne existierenden Kunden. Mit \icode{ON DELETE} und \icode{ON UPDATE} legt man fest, was mit den Child-Datensätzen passiert:

\begin{stdtable}{l X}{Option & Verhalten}
  \icode{RESTRICT} / \icode{NO ACTION} & Löschen bzw. Ändern des Parent-Datensatzes wird verhindert. \\
  \icode{CASCADE} & Child-Datensätze werden mitgelöscht bzw. mitgeändert. \\
  \icode{SET NULL} & Der Fremdschlüssel im Child wird auf \icode{NULL} gesetzt. \\
\end{stdtable}

\subsubsection{Löschen: DROP, TRUNCATE, DELETE}
Mit \icode{DROP} wird eine Tabelle oder Datenbank komplett entfernt, mit \icode{TRUNCATE} werden nur alle Zeilen gelöscht -- die Tabelle selbst bleibt bestehen.

\begin{syntaxbox}[SQL]
DROP DATABASE kurs_db;          -- Datenbank komplett entfernen
DROP TABLE Personal;            -- Tabelle komplett entfernen
TRUNCATE TABLE Personal;        -- alle Zeilen loeschen, Struktur bleibt
\end{syntaxbox}

\begin{stdtable}{l l X}{Befehl & Löscht Daten? & Löscht Struktur?}
  \icode{DROP TABLE} & Ja & Ja -- Tabelle ist weg \\
  \icode{TRUNCATE TABLE} & Ja & Nein -- Tabelle bleibt \\
  \icode{DELETE FROM} & Ja (mit \icode{WHERE} filterbar) & Nein -- Tabelle bleibt \\
\end{stdtable}

\begin{aufgabe}[Tabelle klausurnah anlegen]
Erstelle eine Tabelle \icode{mitarbeiter} mit automatisch nummeriertem Primärschlüssel, eindeutiger Personalnummer und Fremdschlüssel auf \icode{abteilung}.
\begin{syntaxbox}[SQL]
CREATE TABLE mitarbeiter (
  id_mitarbeiter INT PRIMARY KEY AUTO_INCREMENT,
  personalnummer VARCHAR(6) UNIQUE,
  name           VARCHAR(60) NOT NULL,
  id_abteilung   INT,
  FOREIGN KEY (id_abteilung) REFERENCES abteilung(id_abteilung)
);
\end{syntaxbox}
\end{aufgabe}

\subsection{DML -- Data Manipulation Language}
Die DML umfasst alle Befehle, die Daten \emph{innerhalb} bestehender Tabellen verändern: einfügen (\icode{INSERT}), ändern (\icode{UPDATE}), ersetzen (\icode{REPLACE}) und löschen (\icode{DELETE}).

\subsubsection{INSERT}
\begin{syntaxbox}[SQL]
-- Mit Spaltenliste (empfohlen), mehrere Datensaetze auf einmal
INSERT INTO Leistung (LNr, Bezeichnung, Kosten, AID) VALUES
    (800, 'FamilienRecht',  67.00, NULL),
    (801, 'ArbeitsRecht',  700.00, NULL);

-- Ohne Spaltenliste: nur wenn ALLE Spalten in richtiger Reihenfolge angegeben
INSERT INTO Leistung VALUES (804, 'Sonstiges', 99.00, NULL);
\end{syntaxbox}

\begin{infobox}[Hinweis zur Spaltenliste]
Die runden Klammern nach dem Tabellennamen (die Spaltenliste) kann man weglassen, wenn man alle Spalten in der korrekten Reihenfolge angibt. Die Klammern bei \icode{VALUES (...)} bleiben immer.
\end{infobox}

\subsubsection{UPDATE, REPLACE und DELETE}
\begin{syntaxbox}[SQL]
-- Eine oder mehrere Spalten aendern
UPDATE Leistung
SET AID = 2, Kosten = 99.00
WHERE LNr = 17;

-- Bestehenden Datensatz komplett ersetzen
REPLACE kurs_db.Abteilung VALUES (4, 'Rechnungswesen');

-- Datensaetze mit Bedingung loeschen
DELETE FROM Leistung
WHERE LNr = 800;
\end{syntaxbox}

\begin{warnbox}[Achtung]
\icode{UPDATE} und \icode{DELETE} ohne \icode{WHERE} verändern bzw. löschen \emph{alle} Zeilen der Tabelle. Vor dem Ausführen immer die \icode{WHERE}-Bedingung prüfen.
\end{warnbox}

\subsection{DQL -- Data Query Language}
\icode{SELECT} liest Daten aus einer oder mehreren Tabellen. Das Ergebnis ist immer wieder eine Tabelle (Mengeneigenschaft).

\subsubsection{Aufbau eines SELECT-Statements}
Die Schreibreihenfolge ist fest vorgegeben:
\begin{syntaxbox}[SQL]
SELECT   Spalten / Aggregatfunktionen
FROM     Tabellenname
WHERE    Bedingung          -- Zeilen filtern (VOR der Gruppierung)
GROUP BY Spalte             -- Gruppen bilden
HAVING   Bedingung          -- Gruppen filtern (NACH der Gruppierung)
ORDER BY Spalte ASC / DESC  -- Sortieren
LIMIT    Anzahl;            -- Anzahl begrenzen
\end{syntaxbox}

Die Schreibreihenfolge und die interne Ausführungsreihenfolge sind unterschiedlich:

\kzweispaltig{
\begin{stdtable}{l l}{Schritt & Befehl}
  1 & \icode{FROM} \\
  2 & \icode{WHERE} \\
  3 & \icode{GROUP BY} \\
  4 & \icode{HAVING} \\
  5 & \icode{SELECT} \\
  6 & \icode{ORDER BY} \\
  7 & \icode{LIMIT} \\
\end{stdtable}
}{
\begin{stdtable}{l X}{Begriff & Bedeutung}
  Horizontal & bestimmte \emph{Zeilen} filtern (\icode{WHERE}) \\
  Vertikal & bestimmte \emph{Spalten} auswählen (\icode{SELECT s1, s2}) \\
\end{stdtable}
}

\begin{syntaxbox}[SQL]
SELECT * FROM MA;                -- * = alle Spalten
SELECT VName, Gehalt FROM MA;    -- nur bestimmte Spalten (vertikal)
SELECT name AS Kundenname FROM kunde;  -- AS = Spalten-Alias
\end{syntaxbox}

\subsubsection{WHERE -- Zeilen filtern}
Mit \icode{WHERE} definiert man Einschränkungen (horizontales Filtern).

\begin{stdtable}{l X}{Operator & Bedeutung}
  \icode{>} / \icode{>=} & größer / größer gleich \\
  \icode{<} / \icode{<=} & kleiner / kleiner gleich \\
  \icode{=} & Gleichheit \\
  \icode{!=} oder \icode{<>} & ungleich \\
  \icode{AND} & beide Bedingungen müssen wahr sein \\
  \icode{OR} & mindestens eine Bedingung muss wahr sein \\
  \icode{NOT} & kehrt die Bedingung um \\
  \icode{XOR} & genau eine Bedingung darf wahr sein \\
\end{stdtable}

Klammern steuern die Reihenfolge der Auswertung:
\begin{syntaxbox}[SQL]
SELECT id_kunde, name FROM kunde
WHERE (ort = 'Muenster' OR YEAR(geb_datum) > 1960) AND id_kunde > 2;

-- BETWEEN: Bereich (inklusive Grenzen)
WHERE Gehalt BETWEEN 2000 AND 4000
-- entspricht: WHERE Gehalt >= 2000 AND Gehalt <= 4000

-- LIKE: Mustererkennung
WHERE VName LIKE 'M%'   -- beginnt mit M (% = beliebig viele Zeichen)
WHERE VName LIKE 'M_'   -- M + genau 1 Zeichen (_ = genau ein Zeichen)

-- IN: ersetzt eine Kette an OR-Bedingungen
WHERE AID IN (1, 2, 3)
-- entspricht: WHERE AID = 1 OR AID = 2 OR AID = 3
\end{syntaxbox}

\subsubsection{Funktionen}
Aggregatfunktionen fassen mehrere Zeilen zu einem einzelnen Wert zusammen. Sie arbeiten immer auf einer ganzen Spalte oder Gruppe. Daneben gibt es Zeichenketten- und Zeitfunktionen.

\begin{stdtable}{l X}{Funktion & Bedeutung}
  \icode{COUNT(*)} & Anzahl der Zeilen \\
  \icode{SUM(spalte)} & Summe aller Werte \\
  \icode{AVG(spalte)} & Durchschnitt aller Werte (\icode{NULL} wird ignoriert) \\
  \icode{MIN(spalte)} / \icode{MAX(spalte)} & kleinster / größter Wert \\
  \icode{LENGTH(s)} & Anzahl der Zeichen: \icode{LENGTH('frie')} $\rightarrow$ 4 \\
  \icode{LEFT(s, n)} / \icode{RIGHT(s, n)} & die ersten / letzten $n$ Zeichen einer Zeichenkette \\
  \icode{CURDATE()} / \icode{CURTIME()} & aktuelles Datum / aktuelle Uhrzeit \\
  \icode{CURRENT\_TIMESTAMP} & aktuelles Datum samt Uhrzeit \\
  \icode{YEAR(datum)} & Jahr eines Datums \\
\end{stdtable}

\begin{syntaxbox}[SQL]
SELECT
    COUNT(*)    AS AnzahlMitarbeiter,  -- AS vergibt einen Alias
    AVG(Gehalt) AS Durchschnittsgehalt,
    MIN(Gehalt) AS MinGehalt,
    MAX(Gehalt) AS MaxGehalt,
    SUM(Gehalt) AS Gesamtlohn
FROM MA;
\end{syntaxbox}

\subsubsection{GROUP BY und HAVING}
\icode{GROUP BY} fasst Zeilen mit dem gleichen Wert in einer Spalte zu einer Gruppe zusammen. Danach kann man pro Gruppe eine Aggregatfunktion anwenden. \icode{WHERE} filtert Zeilen \emph{vor} der Gruppierung, \icode{HAVING} filtert Gruppen \emph{nach} der Gruppierung.

\begin{syntaxbox}[SQL]
-- Ohne GROUP BY: ein Wert ueber alle Zeilen
SELECT AVG(Gehalt) FROM MA;

-- Mit GROUP BY: ein Wert pro Abteilung
SELECT AID, AVG(Gehalt) AS Durchschnittsgehalt
FROM MA
GROUP BY AID;

SELECT AID, COUNT(*) AS Anzahl
FROM MA
WHERE Gehalt > 2000       -- filtert Zeilen BEVOR Gruppen gebildet werden
GROUP BY AID
HAVING COUNT(*) > 2;      -- filtert Gruppen NACHDEM sie gebildet wurden
\end{syntaxbox}

\begin{merksatz}[Wichtige Regel]
Jede Spalte im \icode{SELECT}, die keine Aggregatfunktion ist, muss auch im \icode{GROUP BY} stehen -- sonst gibt es einen Fehler.
\end{merksatz}

\subsubsection{ORDER BY und LIMIT}
\begin{syntaxbox}[SQL]
SELECT VName, Gehalt FROM MA
ORDER BY Gehalt DESC;   -- DESC = absteigend, ASC = aufsteigend (Standard)

SELECT * FROM MA LIMIT 5;                  -- MySQL
SELECT * FROM MA FETCH FIRST 5 ROWS ONLY;  -- Standard-SQL
\end{syntaxbox}

\subsubsection{Joins}
Joins verknüpfen Daten aus mehreren Tabellen über eine Beziehung -- meist Primärschlüssel (Parent) = Fremdschlüssel (Child).

\begin{stdtable}{l X}{Join-Typ & Was wird ausgegeben?}
  \icode{INNER JOIN} & Nur Datensätze mit Treffer in \emph{beiden} Tabellen (Schnittmenge). \\
  \icode{LEFT JOIN} & Alle Datensätze der linken Tabelle -- auch ohne Treffer rechts. \\
  \icode{RIGHT JOIN} & Alle Datensätze der rechten Tabelle -- auch ohne Treffer links. \\
  \icode{CROSS JOIN} & Kartesisches Produkt: jede Zeile wird mit jeder kombiniert. \\
\end{stdtable}

\begin{syntaxbox}[SQL]
SELECT bestelldatum, name FROM kunde
INNER JOIN bestellung ON kunde.id_kunde = bestellung.id_kunde;
\end{syntaxbox}

\begin{infobox}[Merkhilfe: Schnittmengendiagramm]
\icode{INNER} = Überschneidung der beiden Kreise, \icode{LEFT}/\icode{RIGHT} = der ganze linke bzw. rechte Kreis, \icode{CROSS} = alle Kombinationen.
\end{infobox}

\subsubsection{Unterabfragen (Subselects)}
Wie in der Mathematik bei $f(g(x))$: Das innere Statement wird zuerst ausgeführt, sein Ergebnis nutzt die äußere Abfrage weiter. Das Ergebnis einer Unterabfrage ist immer eine Tabelle. Vorteil: Ein Teilproblem wird getrennt gelöst -- oft eine übersichtlichere Alternative zum Join.

\begin{syntaxbox}[SQL]
-- Unterabfrage im FROM (braucht immer einen Alias)
SELECT * FROM
    (SELECT VName FROM MA) AS temp;  -- AS temp ist Pflicht

-- Subselect im WHERE: zuerst MAX berechnen, dann damit vergleichen
SELECT VName, Gehalt FROM MA
WHERE Gehalt = (SELECT MAX(Gehalt) FROM MA);
\end{syntaxbox}

\subsubsection{Views (Sichten)}
Eine \emph{View} ist eine gespeicherte \icode{SELECT}-Abfrage, die wie eine virtuelle Tabelle genutzt wird. Im Unterschied zu einer Tabelle speichert eine View keine eigenen Daten, sondern nur die Abfrage -- die Daten liegen weiter in den Basistabellen.

\begin{procon}
  \pro{Komplexe Abfragen unter einem Namen speichern}
  \pro{Einheitliche Darstellung für alle Nutzer}
  \pro{Eingeschränkter, datensparsamer Zugriff}
  \con{Keine eigenen Daten -- Performance hängt an den Basistabellen}
\end{procon}

\subsection{DCL -- Data Control Language}
Die DCL steuert Zugriff und Rechte auf die Datenbank. Kernfrage des Datenschutzes: \emph{Wer darf welche Daten zu welchem Zweck sehen oder ändern?}

\begin{merksatz}[Grundsatz]
Jeder User bekommt nur die Rechte, die er wirklich benötigt -- kein pauschales \icode{ALL}.
\end{merksatz}

\icode{GRANT} vergibt Rechte, \icode{REVOKE} entzieht sie. Rechte sind u.\,a. \icode{ALL}, \icode{SELECT}, \icode{INSERT}, \icode{UPDATE}, \icode{DELETE}, \icode{DROP}. Es gibt drei Ebenen: global (\icode{*.*}), Datenbank (\icode{db.*}) und Tabelle (\icode{db.tabelle}).

\begin{syntaxbox}[SQL]
-- Nutzer anlegen und loeschen
CREATE USER 'test'@'localhost' IDENTIFIED BY 'pw123';
DROP USER   'test'@'localhost';

-- Rechte vergeben und entziehen
GRANT SELECT ON kurs_db.kunde TO 'test'@'localhost';
GRANT ALL    ON kurs_db.*     TO 'test'@'localhost';
REVOKE ALL   ON kurs_db.mitarbeiter FROM 'test'@'localhost';
\end{syntaxbox}

\subsection{TCL -- Transaktionen}
Eine \emph{Transaktion} ist eine logisch zusammenhängende Gruppe von SQL-Befehlen, die vollständig oder gar nicht ausgeführt wird -- z.\,B. eine Überweisung: abbuchen \emph{und} gutschreiben.

\begin{syntaxbox}[SQL]
START TRANSACTION;
UPDATE konto SET saldo = saldo - 100 WHERE id_konto = 1;
UPDATE konto SET saldo = saldo + 100 WHERE id_konto = 2;
COMMIT;     -- dauerhaft speichern
-- ROLLBACK; -- alternativ: alle Aenderungen verwerfen
\end{syntaxbox}

Transaktionen arbeiten zuverlässig nach dem \icode{ACID}-Prinzip:
\begin{stdtable}{l X}{Eigenschaft & Bedeutung}
  \icode{A}tomicity & Alles oder nichts -- eine Transaktion wird nie teilweise ausgeführt. \\
  \icode{C}onsistency & Die Datenbank ist vor und nach der Transaktion konsistent. \\
  \icode{I}solation & Gleichzeitige Transaktionen beeinflussen sich nicht. \\
  \icode{D}urability & Nach dem \icode{COMMIT} sind die Änderungen dauerhaft gespeichert. \\
\end{stdtable}

\subsubsection{Locking}
Beim \emph{Locking} werden Datenbankobjekte (Zelle, Zeile, Tabelle oder ganze Datenbank) temporär gesperrt; die Freigabe erfolgt meist nach dem Transaktionsende.
\begin{itemize}
  \item \textbf{Shared Lock (Read Lock):} Mehrere Transaktionen dürfen gleichzeitig lesen, keine darf schreiben.
  \item \textbf{Exclusive Lock (Write Lock):} Nur eine Transaktion darf lesen und schreiben; alle anderen müssen warten.
\end{itemize}
Gegenseitiges Sperren kann zu einem \emph{Deadlock} führen.

\subsubsection{Isolationsebenen}
Isolationsebenen legen fest, wie stark gleichzeitige Transaktionen voneinander getrennt werden (von der schwächsten zur stärksten): \icode{READ UNCOMMITTED} $\rightarrow$ \icode{READ COMMITTED} $\rightarrow$ \icode{REPEATABLE READ} $\rightarrow$ \icode{SERIALIZABLE}. Höhere Isolation bedeutet weniger Nebeneffekte, aber geringere Parallelität.

\subsubsection{Events}
Ein \emph{Event} ist eine automatisch geplante SQL-Anweisung, die zu einer bestimmten Zeit oder regelmäßig ausgeführt wird -- z.\,B. um alte Daten zu löschen, Auswertungen vorzubereiten oder Logs zu bereinigen.

\subsection{Trigger, Funktionen und Prozeduren}
\subsubsection{Trigger}
Ein \emph{Trigger} fängt ein Ereignis ab (\icode{INSERT}, \icode{UPDATE} oder \icode{DELETE} an einer Tabelle) und führt automatisch einen Anweisungsblock aus („feuern"). Er wird nie vom Benutzer aufgerufen, sondern vom DBMS -- entweder \icode{BEFORE} oder \icode{AFTER} der Änderung. Typische Einsatzzwecke: Werte automatisch setzen, Plausibilität prüfen, Änderungen protokollieren und referenzielle Integrität überwachen.

\subsubsection{Funktionen}
Eine benutzerdefinierte Funktion ist eine \emph{Skalarfunktion}: Sie gibt genau einen Wert zurück. Der Rückgabetyp wird mit \icode{RETURNS} deklariert, der Wert mit \icode{RETURN} zurückgegeben. Nach \icode{DETERMINISTIC} folgt der eigentliche Funktionskörper.

\begin{syntaxbox}[SQL]
CREATE FUNCTION multipliziere(param1 FLOAT(8,2), param2 FLOAT(8,2))
RETURNS FLOAT(8,2)
READS SQL DATA
DETERMINISTIC
RETURN param1 * param2;

SELECT multipliziere(Kosten, 1.19) FROM Leistung;   -- Aufruf
\end{syntaxbox}

\subsubsection{Prozeduren (Stored Procedures)}
Eine \emph{Stored Procedure} ist ein benannter, in der Datenbank gespeicherter Anweisungsblock -- mit oder ohne Parameter, aufgerufen per \icode{CALL}. Nutzen: Businesslogik in die Datenbank auslagern, Wiederverwendbarkeit von mehreren Rechnern aus und Parameterübergabe.

\newpage
